% Vergleich zur ausdrücklich geometrischen Zwölferteilung.
% Unveränderte geometrische Sollwerte: A4, 17,5 / 35 / 24,75 / 49,5 mm.
% typearea passt die Texthöhe an ganze Zeilen an; diese Datei zeigt das Ergebnis.
\documentclass[12pt,a4paper,ngerman,twoside,DIV=12,BCOR=0mm]{scrartcl}
\usepackage{iftex}
\usepackage[T1]{fontenc}
\ifPDFTeX
  \usepackage[utf8]{inputenc}
\else
  \usepackage[utf8]{luainputenc}
\fi
\usepackage[type1]{ebgaramond}
\usepackage[type1]{FiraSans}
\usepackage[ngerman]{babel}
\usepackage{microtype}
\usepackage{scrlayer-scrpage}
\setkomafont{disposition}{\rmfamily\bfseries}
\setkomafont{section}{\rmfamily\normalsize\bfseries}
\setkomafont{pageheadfoot}{\sffamily\small\firalining\firatabular}
\recalctypearea
\ExplSyntaxOn
\NewDocumentCommand{\inmm}{m}{\fp_eval:n{round(\dim_to_fp:n{#1}/\dim_to_fp:n{1mm},3)}\,mm}
\ExplSyntaxOff
\begin{document}
\section*{Zwölferteilung mit typearea}
Diese Datei verwendet ausschließlich \textsf{DIV=12,BCOR=0mm}.
Anders als in den geometrischen Beispielen wird keine zusätzliche
\textsf{geometry}-Vorgabe angewendet.

Nach dem Laden der Schrift wird der Satzspiegel neu berechnet.
Die Höhe richtet sich dabei nach den Zeilen des eingestellten Schriftgrads.
Gleiche Proportionen bedeuten daher nicht automatisch exakt gleiche Maße.

{\sffamily\firalining\firatabular
\begin{tabular}{ll}
Papierbreite & \inmm{\paperwidth}\\
Papierhöhe & \inmm{\paperheight}\\
Textbreite & \inmm{\textwidth}\\
Texthöhe & \inmm{\textheight}\\
Innenrand & \inmm{1in+\oddsidemargin}\\
Oberrand bis zum Text & \inmm{1in+\topmargin+\headheight+\headsep}
\end{tabular}}

Die geometrischen Vergleichswerte betragen 157,5\,mm Textbreite
und 222,75\,mm Texthöhe; innen 17,5\,mm und oben 24,75\,mm.
\end{document}
